%% kapa: exam=1 round=28 year=2017 period=1 form=A subject=real_estate subject_ko=부동산학원론 seq=4 number=84
%% source: Q-Net 감정평가사 기출문제 (cst00309) artlSeq=5199603 file=extracted/1008/5199603/2017년 제28회 감정평가사 1차시험 문제지/1교시 A형.pdf
%% conversion: kordoc 4.24.51 -> tools/gen1.py (auto); review=proofread-sonnet; flags=none
\begin{question}{84}
A금융기관은 원금균등분할상환방식과 원리금균등분할상환방식의 대출을 제공하고 있다. 두 방식에 의해 산정한 첫번째 월불입액의 차액은? (단, 주어진 조건에 한함)
\begin{framedtext}
○ 주택가격: 6억원 \par
○ 담보인정비율(LTV): 50\% \par
○ 대출조건(매월말 상환): 대출기간은 30년, 대출이자율은 연 6\%(월 0.5\%, 월 저당상수 = 0.006443) \par
○ 원금균등분할상환방식: 3년 거치 후 원금균등분할상환하며, 거치기간 동안에는 이자만 지급함 \par
○ 원리금균등분할상환방식: 거치기간 없음
\end{framedtext}
\begin{choices}
  \item 332,900원
  \item 432,900원
  \item 532,900원
  \item 632,900원
  \item 732,900원
\end{choices}
\end{question}
\answer{2}
